\chapter{Reference: composition}
\label{ch:refman-compose}

\begin{aside}
Big apps are made of small views. The trick is to write each
view as a function and call it from the parent. State stays
in one Model; rendering decomposes.
\end{aside}

\section{Working example}

\begin{lstlisting}[language=C++]
#include <maya/maya.hpp>

using namespace maya;
using namespace maya::dsl;

struct Compose {
    struct Model {
        std::string title = "dashboard";
        int cpu = 42;
        int mem = 67;
    };
    struct Quit {};
    using Msg = std::variant<Quit>;

    static Model init() { return {}; }
    static auto update(Model m, Msg)
        -> std::pair<Model, Cmd<Msg>> {
        return {m, Cmd<Msg>::quit()};
    }

    // -- Small view helpers --
    static Element header(const Model& m) {
        return h(
            text(m.title) | Bold,
            spacer_,
            t<"q to quit"> | Dim
        );
    }

    static Element gauge(std::string label, int pct) {
        std::string bar;
        for (int i = 0; i < 20; ++i) {
            bar += (i < pct / 5) ? "#" : ".";
        }
        return h(
            text(label) | Bold,
            text(": "),
            text(bar) | Fg<100, 220, 160>,
            text(std::format(" {}\%", pct))
        );
    }

    static Element body(const Model& m) {
        return v(
            gauge("cpu", m.cpu),
            gauge("mem", m.mem)
        );
    }

    // -- Top-level view glues them together --
    static Element view(const Model& m) {
        return v(
            header(m),
            blank_,
            body(m)
        ) | pad<1> | border_<Round>;
    }

    static auto subscribe(const Model&) -> Sub<Msg> {
        return key_map<Msg>({ {'q', Quit{}} });
    }
};

int main() { run<Compose>({.title = "compose"}); }
\end{lstlisting}

\section{Notes}

\begin{itemize}[leftmargin=*]
  \item Pass the Model by const reference into helpers;
    don't reach for globals.
  \item Helpers that take just the data they need are
    easier to test and refactor.
  \item For very large apps, split helpers across files; the
    \code{Element} return type travels fine across TUs.
\end{itemize}

\section{Recap}

\begin{recap}
\begin{itemize}[leftmargin=*]
  \item Helpers return \code{Element}.
  \item Pass only what they need.
  \item Top-level \code{view} glues them.
\end{itemize}
\end{recap}

\section*{Workshop}

\begin{enumerate}[leftmargin=*]
  \item Extract three helpers from your current view.
  \item Move them to a separate \code{views.cpp}.
  \item Pass a parameter to make one helper reusable.
\end{enumerate}
